% !TEX encoding = UTF-8 Unicode
\chapter{Sintaxis Híbrida Lisp/Python}
\label{chap:hybrid-syntax}
{\huge$\bigstar$}{\huge$\bigstar$}{\huge$\bigstar$} Avanzado

LisPy permite escribir archivos \texttt{.lp} que mezclan Lisp y Python en un mismo
archivo, sin necesidad de \texttt{py:run} explícito. El preprocesador híbrido separa
automáticamente los bloques Lisp (líneas que empiezan con \texttt{(}) de los bloques Python.

\section{Mecanismo del Preprocesador}

El preprocesador opera en tres fases:
\begin{enumerate}
\item \textbf{Detección de bloques}: Líneas que comienzan con \texttt{(} son Lisp;
      el resto es Python. Bloques Python se agrupan por indentación.
\item \textbf{Resolución de símbolos}: Los nombres de funciones Lisp usados en
      Python se reemplazan por \texttt{builtins.nombre\_cb(args)}.
      Las referencias \texttt{@variable} en Python se convierten en
      \texttt{builtins.lispy\_callback\_c('variable')}.
\item \textbf{Inyección}: Se generan automáticamente \texttt{py:register-callback}
      para cada \texttt{defun} usado en Python, y cada bloque Python se envuelve
      en \texttt{py:run}.
\end{enumerate}

\section{Ejemplo Básico}

\begin{lispcode}
;;; Archivo hibrido.lp
(defun sumar (a b) (+ a b))
(defvar *factor* 10)

x = sumar(3, 4)
resultado = x * factor

(print @resultado)
\end{lispcode}

El preprocesador lo convierte automáticamente en:

\begin{lispcode}
(defun sumar (a b) (+ a b))
(defvar *factor* 10)
(py:register-callback 'sumar)
(py:run "
import builtins
x = builtins.sumar_cb(3, 4)
resultado = x * builtins.lispy_callback_c('*factor*')
")
(print @resultado)
\end{lispcode}

\section{Carga de Archivos Híbridos}

\begin{lispcode}
;;; Carga explícita
(py:load-hybrid "archivo.lp")

;;; Carga automática: load detecta archivos híbridos
(load "archivo.lp")
\end{lispcode}

\section{Tabla de Ambigüedades}

\begin{center}
\begin{longtable}{ll}
\toprule
\textbf{Caso} & \textbf{Resolución} \\
\midrule
\texttt{@var} en Lisp & \texttt{py:\%chain "var"} (sin cambios) \\
\texttt{@var} en Python & \texttt{builtins.lispy\_callback\_c('var')} \\
\texttt{fn(args)} (fn es Lisp) & \texttt{builtins.fn\_cb(args)} \\
\texttt{print(...)} en Python & Se mantiene como Python \\
\texttt{x = (lisp-form)} & Error: no mezclar en misma línea \\
\bottomrule
\end{longtable}
\end{center}

\section{Ejemplos Avanzados}

\subsection{Funciones Lisp en F-strings de Python}
\begin{lispcode}
(defun doble (n) (* 2 n))

x = 21
print(f"El doble de {x} es {doble(x)}")
\end{lispcode}

\subsection{Iteración Cruzada}
\begin{lispcode}
(defun factorial (n)
  (if (<= n 1) 1 (* n (factorial (- n 1)))))

for i in range(1, 11):
    print(f"{i}! = {factorial(i)}")
\end{lispcode}

\subsection{Variables Lisp en Python}
\begin{lispcode}
(defparameter *pi* 3.1415926535)
(defvar *radio* 5)

area = pi * radio ** 2
print(f"Área del círculo: {area}")
\end{lispcode}

\section{Limitaciones}

\begin{itemize}
\item No se permite mezclar Lisp y Python en la misma línea.
\item Las funciones Lisp usadas en Python deben estar definidas antes del bloque Python.
\item Los callbacks son sincrónicos: cada llamada Python$\to$Lisp cruza el puente FFI.
\end{itemize}

\begin{notebox}[title=Rendimiento]
Para bucles intensivos (más de 1000 iteraciones con callbacks Lisp), considera
mover la lógica completa a Lisp o a Python puro con \texttt{py:run}. El overhead
del puente es significativo en llamadas frecuentes.
\end{notebox}
